\documentclass{article}
\usepackage[utf8]{inputenc}
\usepackage{booktabs}
\usepackage{array}
\usepackage{geometry}
\geometry{a4paper, margin=1in}
\usepackage[backend=biber, style=ieee]{biblatex}
\addbibresource{bibliography.bib}

\begin{document}

\section*{Literature Review}

\subsection*{Detailed Review of Core Studies}

\begin{enumerate}
    \item Not All Languages Are Created Equal in LLMs: Improving Multilingual Capability by Cross-Lingual-Thought Prompting \textbf{\cite{huang-etal-2023-languages}}
    \begin{itemize}
        \item \textbf{What it does:} Introduces Cross-Lingual-Thought (XLT) prompting, a strategy that encourages Large Language Models (LLMs) to engage in cross-lingual and logical reasoning by translating or rephrasing prompt queries into English before executing step-by-step reasoning.
        \item \textbf{On what data:} Evaluated across arithmetic reasoning, multi-hop reasoning, and multilingual NLP benchmarks in multiple languages (e.g., MGSM, Native Reasoning datasets).
        \item \textbf{With what result:} Demonstrates consistent performance gains over direct native-language prompting by bridging cross-lingual reasoning gaps and leveraging English-centric knowledge representation.
    \end{itemize}

    \item Is Translation All You Need? A Study on Solving Multilingual Tasks \textbf{\cite{liu-etal-2025-translation}}
    \begin{itemize}
        \item \textbf{What it does:} Investigates the trade-offs and limitations of translation-based prompting versus native-language prompting, evaluating when English translation helps or hinders task performance.
        \item \textbf{On what data:} Evaluated on standard NLP reasoning tasks as well as culture-specific and socio-linguistic benchmarks across high- and low-resource languages.
        \item \textbf{With what result:} Finds that while translating to English improves performance on standard NLP tasks for English-centric models, native-language prompting significantly outperforms translation on culture-related tasks that require deep local context.
    \end{itemize}

    \item Do Llamas Work in English? On the Latent Language of Multilingual Transformers \textbf{\cite{wendler-etal-2024-llamas}}
    \begin{itemize}
        \item \textbf{What it does:} Applies mechanistic interpretability techniques (specifically the logit lens) to analyze intermediate layer representations of multilingual transformer models (e.g., Llama-2 family) during non-English processing.
        \item \textbf{On what data:} Tested on parallel multilingual text inputs across high-, medium-, and low-resource languages.
        \item \textbf{With what result:} Maps model internal transitions showing a progression from the input language space to an intermediate, English-biased ``concept space,'' and finally decoding into the target language output space.
    \end{itemize}

    \item MMLU-ProX: A Multilingual Benchmark for Advanced Large Language Model Evaluation \textbf{\cite{xuan-etal-2025-mmlu}}
    \begin{itemize}
        \item \textbf{What it does:} Introduces MMLU-ProX, a massively multilingual benchmark designed for robust evaluation of reasoning and knowledge across typologically diverse languages using parallel, high-quality questions.
        \item \textbf{On what data:} Benchmark suite comprising parallel evaluation sets across 29 diverse languages spanning multiple domains and subjects.
        \item \textbf{With what result:} Reveals severe performance degradation in state-of-the-art LLMs on low-resource and non-Western languages (e.g., African languages) compared to English and high-resource languages.
    \end{itemize}

    \item Measuring Massive Multitask Language Understanding \textbf{\cite{hendrycks2021measuring}}
    \begin{itemize}
        \item \textbf{What it does:} Proposes a new massive multitask language understanding test covering 57 tasks (including STEM, humanities, social sciences, law, and ethics) evaluated under zero-shot and few-shot settings to measure a model's extensive world knowledge and problem-solving abilities.
        \item \textbf{On what data:} 15,908 total multiple-choice questions collected from freely available online sources, professional practice exams, and university-level materials, split into few-shot development, validation, and test sets.
        \item \textbf{With what result:} Finds that while the 175B parameter GPT-3 model outperforms random chance (reaching 43.9\% average accuracy), performance remains subhuman across all tasks, exhibits lopsided weaknesses in calculation-heavy and value-laden subjects (like law and morality), and shows significant model miscalibration.
    \end{itemize}

    \item Language Model Tokenizers Introduce Unfairness Between Languages \textbf{\cite{petrov2023language}}
    \begin{itemize}
        \item \textbf{What it does:} Demonstrates how disparities in the treatment of different languages arise systematically at the tokenization stage prior to model invocation, introducing the concept of "tokenizer parity" and advocating for multilingually fair subword tokenizers.
        \item \textbf{On what data:} Evaluated across the FLORES-200 parallel corpus spanning 200 languages for various English-centric, non-English centric, multilingual, and byte-level tokenizers.
        \item \textbf{With what result:} Shows that tokenization lengths for identical semantic content can vary by up to 15 times (e.g., for Shan compared to English), directly introducing severe economic, latency, and context-length inequalities.
    \end{itemize}
\end{enumerate}

\clearpage

\subsection*{Comparative Summary of Published Numbers}

\begin{table}[htbp]
\centering
\caption{Comparison of published evaluation metrics across models, datasets, and specific splits.}
\label{tab:published_numbers}
\begin{tabular}{p{4.5cm} p{3.5cm} p{2.5cm} p{4.5cm}}
\toprule
\textbf{Study} & \textbf{Dataset} & \textbf{Split / Subset} & \textbf{Published Result / Metric} \\
\midrule
Huang et al. (2023) & MGSM (Multilingual Grade School Math) & Multilingual Test Split (XLT Prompting) & $+8.5\%$ average accuracy improvement over Direct Prompting \\
\midrule
Liu et al. (2025) & CulturalNLP / Standard NLP & Culture-Specific Test Split vs. Standard Split & Native Prompting outperforming English Translation by up to $12.4\%$ on cultural tasks \\
\midrule
Wendler et al. (2024) & Llama-2 Internal Latent Logit Lens Evaluation & Intermediate Layers (Layers 12--24) & $>70\%$ probability alignment with English concept space prior to target language decoding \\
\midrule
Xuan et al. (2025) & MMLU-ProX & High-Resource Split (e.g., FR/DE/ES) & Average Accuracy: $68.2\%$ \\
 & MMLU-ProX & Low-Resource Split (e.g., African Languages) & Average Accuracy: $31.4\%$ (Performance drop of $\sim36.8\%$) \\
\midrule
Hendrycks et al. (2021) & MMLU (Massive Multitask Test) & GPT-3 X-Large (175B) Few-Shot Test Set & Average Accuracy: $43.9\%$ (ranging from $69\%$ on US Foreign Policy down to $26\%$ on College Chemistry)\\
\midrule
Petrov et al. (2023) & FLORES-200 Parallel Corpus & ChatGPT / GPT-4 Tokenizer ($\text{cl100k\_base}$) & Tokenization premiums up to $\sim15\times$ longer for low-resource languages (e.g., Shan) relative to English \\
\bottomrule
\end{tabular}
\end{table}
\section{Proposed Metrics}
To measure the cross-lingual gap, we will cover two primary quantitative dimensions corresponding to response quality and token efficiency, each measured by an appropriate metric:

\subsubsection{Quality Metric: Cross-Lingual Parity Ratio ($R_{\text{parity}}$)}
To measure the \textbf{quality gap} across MMLU, GSM8K, ARC-Challenge, and BoolQ, we use the Cross-Lingual Parity Ratio. This normalizes model performance on parallel Polish benchmark splits ($\mathcal{D}_{\text{PL}}$) against English benchmark splits ($\mathcal{D}_{\text{EN}}$):

\begin{equation}
R_{\text{parity}}(M, \mathcal{D}) = \frac{\text{Score}_{\text{Polish}}(M, \mathcal{D}_{\text{PL}})}{\text{Score}_{\text{English}}(M, \mathcal{D}_{\text{EN}})}
\end{equation}

where $\text{Score}$ represents task accuracy or exact-match. An $R_{\text{parity}}$ close to $1.0$ indicates minimal quality degradation, whereas lower values quantify the relative reasoning loss in Polish.

\subsubsectionEfficiency Metric: Token Efficiency Ratio ($E_{\text{token}}$)}
To measure the \textbf{token efficiency gap}, we evaluate the computational and context-window cost of processing and generating in Polish versus English. Because tokenizers often exhibit high fragmentation for Slavic languages, we define the Token Efficiency Ratio as the ratio of tokens consumed to achieve equivalent correct outputs:

\begin{equation}
E_{\text{token}}(M, \mathcal{D}) = \frac{\text{Tokens}_{\text{Polish}}(M, \mathcal{D}_{\text{PL, correct}})}{\text{Tokens}_{\text{English}}(M, \mathcal{D}_{\text{EN, correct}})}
\end{equation}

This metric captures token inflation (e.g., $E_{\text{token}} > 1.3$), exposing how much extra computational overhead and context budget Polish inference requires compared to English.
\printbibliography

\end{document}